\section{Tight capacities and two scales}\label{sec:scales}

We first reduce the input to one in which every coordinate bound is
attained.  The use of tight capacities and a maximum-capacity spanning
forest follows the representation of integral flows in~\cite{CDR2010}.
The estimates below allow both dimensions and all capacities to vary.

For feasibility, form a directed network with a source, the row and
column vertices, and a sink. Give the source-to-row arcs capacities
$r_i$, the row-to-column arcs capacities $b_{ij}$, and the
column-to-sink arcs capacities $c_j$. An integral flow of value
$N=\sum_i r_i=\sum_j c_j$ is exactly a feasible table: at that value
every arc incident to the source or sink is saturated. The
shortest-augmenting-path algorithm of Edmonds and
Karp~\cite[Theorem~1]{EK1972} uses $O(VE)$ augmentations, where $V$
and $E$ are the numbers of vertices and arcs.  The present network is
simple, so $E=O(V^2)$ and the bound is $O(V^3)$, independently of the
numerical capacities. Each augmenting path is found by breadth-first
search in $O(V+E)$ operations. Starting at zero, augmentation by the
minimum residual capacity preserves integrality. All flow and residual-capacity
values have polynomial binary length, so this gives a deterministic
polynomial-bit feasibility test for the present network.

Return zero if the input is infeasible. For a feasible input, let
$\ell_s$ and $u_s$ be the minimum
and maximum feasible integer values of cell $s$.  Each extremum is
computable by binary search using flow feasibility with one additional
upper or lower bound.  A lower bound is subtracted from its cell and
the two incident margins before applying the flow test; if a residual
margin or capacity is negative, reject that tested bound immediately.
The searches
use polynomially many operations on integers of polynomial binary
length.

Replace each cell value $X_s$ by $X_s-\ell_s$, replace its upper bound
by $b_s=u_s-\ell_s$, and subtract all incident minima from each margin.
This is a bijection between the original and the new feasible sets:
every original feasible table satisfies all its individual extremum
bounds, and adding back the minima is the inverse map.  In the new
instance, each endpoint $0$ and $b_s$ is attained by a feasible table,
possibly a different table for each endpoint and cell.  Delete cells
with $b_s=0$.  A deleted cell contributes no choice.  Retain all row
and column vertices, including isolated vertices, and write $r,c$ for
the new margins and $Z$ for the unchanged count.

Let $L$ be the binary input length from Theorem~\ref{thm:main}, and
retain the original dimensions $m,n$.  Fix
\begin{equation}\label{eq:parameters}
\begin{gathered}
d=mn+m+n+2,\qquad
K=20\bigl(L+d+\lceil\log_2(1/\epsilon)\rceil+1\bigr),
\qquad \gamma=\frac{\epsilon}{512dK},\\
A=\left\lceil\frac{4K}{\gamma}\right\rceil,\qquad
B=16Ad^2,\qquad W=100dB,\qquad H=10Ad^2(W+1).
\end{gathered}
\end{equation}
All these parameters except $\gamma$ are positive integers.
Here $K$ bounds the logarithmic counting scales, $\gamma$ is a relative
width near a tree boundary, and $A$ sets the penalty strength. The
parameter $B$ is the number of bins per free coordinate, $W$ separates
small and large capacities, and $H$ bounds the penalty and is the final
annealing index. Binary encoding and the fact that tightened capacities
never increase give
\begin{equation}\label{eq:capacity-product}
\prod_{s\text{ remaining}}(b_s+1)\le 2^L\le 2^K,
\qquad 2^{-3K}\le\frac{\epsilon}{128}.
\end{equation}
Indeed, a nonnegative integer written with $\ell$ binary digits has
at most $2^\ell$ possible values between zero and itself.  For the
second inequality, the definition of $K$ gives
$3K\ge\log_2(128/\epsilon)$.
Call a remaining cell \emph{small} if $b_s<W$, and \emph{large}
otherwise.  Write $\mathcal S$ for the small cells and set
\[
p=\sum_{s\in\mathcal S}b_s.
\]
There are at most $mn<d$ remaining cells and $m+n<d$ vertices,
so $p\le dW$.  In particular the number of binary slots introduced
later is polynomial.  More explicitly, $d,K=O(L)$ and
\[
A=O(L^3\epsilon^{-1}),\quad
B=O(L^5\epsilon^{-1}),\quad
W=O(L^6\epsilon^{-1}),\quad
H=O(L^{11}\epsilon^{-2}).
\]
Here $L$ includes the encoding of $\epsilon$, so
$\lceil\log_2(1/\epsilon)\rceil=O(L)$.  These bounds concern the
numerical values of the auxiliary parameters; the retained capacities
themselves continue to have polynomial binary length.

Form the bipartite graph whose vertices are the rows and columns and
whose edges are the large cells.  In each connected component choose
a spanning tree of maximum total capacity, and choose a root.  An
isolated vertex constitutes a component with an empty tree.  The
union of these trees is denoted by $\mathcal T$; the large cells
outside $\mathcal T$ form the set $\mathcal F$ of \emph{free cells}.
Sorting large cells by decreasing capacity and adding an edge whenever
it joins two current components constructs such a forest in
polynomial time.

The small cells will have separate row and column coordinates.  The small-cell
contributions to row and column equations will therefore use disjoint
coordinate sets.  We recover an actual table by requiring the two
physical values of each small cell to agree.

For each small cell $s$, take integer counts
$x_s,y_s\in\{0,\ldots,b_s\}$.  Its contribution to its row margin is
$x_s$, and its contribution to its column margin is $b_s-y_s$.
Thus the two contributions agree precisely when $x_s+y_s=b_s$.
For example, when $b_s=3$, the profile $(x_s,y_s)=(2,1)$ gives
physical value two at both endpoints, whereas $(2,2)$ gives row value
two and column value one.
For a large cell there is a single value at both endpoints.  Given
free values $(z_e)_{e\in\mathcal F}$, impose the prescribed margin at
every nonroot vertex and solve for the values $(z_t)_{t\in\mathcal T}$.
As proved below, this has a unique real solution, which is integral
when the free values are integral.  For a root $v$, let $\delta_v$
be its resulting actual incident sum minus its prescribed margin.
Define
\begin{equation}\label{eq:penalty}
D(x,y,z)=A\left(
\sum_{v\text{ a root}}|\delta_v|
+\sum_{t\in\mathcal T}
\frac{\dist(z_t,[0,b_t])}{b_t}
\right).
\end{equation}
In this expression $z$ denotes the free vector; the tree values and
root deviations are its determined functions.  The denominators
$b_t$ are positive because every tree cell is large.

\begin{lemma}[Forest estimates]\label{lem:forest}
For a free cell $e$, every tree cell $t$ on the path joining its
endpoints satisfies $b_e\le b_t$.  For every box-valid integer small profile
$(x,y)$ and every real free vector
$z\in\prod_{e\in\mathcal F}[0,b_e+1]$, the nonroot equations have a
unique solution.  This solution is integral for integral $z$, and
the root deviations are independent of $z$.

For any fixed feasible table $X^*$ of the tightened instance,
\begin{equation}\label{eq:forest-displacement}
|z_t-X^*_t|\le4db_t\quad(t\in\mathcal T),
\qquad
\sum_{v\text{ a root}}|\delta_v|\le2dW.
\end{equation}
Consequently, throughout the extended free box,
\[
0\le D\le H,\qquad
z_t\in[-10db_t,10db_t]\quad(t\in\mathcal T).
\]
For fixed $(x,y)$, set $u_e=z_e/(b_e+1)$.  As a function of $u$,
the penalty $D$ is convex on $[0,1]^{\mathcal F}$ and satisfies
\begin{equation}\label{eq:forest-lipschitz}
|D(x,y,u)-D(x,y,u')|
\le2Ad^2\|u-u'\|_\infty.
\end{equation}
For an empty free set the domain is a singleton and the last bound
has its usual vacuous interpretation.
\end{lemma}

\begin{proof}
If a free edge had larger capacity than a tree edge on its path,
exchanging those two edges would increase the total tree capacity.
This proves the path inequality.  Leaf elimination proves existence
and uniqueness of the nonroot solution: at a nonroot leaf its margin
equation determines the edge to its parent, after which that vertex
can be removed.  The same procedure proves integrality and that all
tree values are affine functions of the free values.  No equation is
imposed at the last, root vertex.  For a singleton component there
are no unknown tree values and no nonroot equations.

Changing a free value affects tree values along its tree path.
Small-cell profiles need more care: their two physical views need not
agree, so their contributions must be counted by incidences rather than
cancelled as an ordinary flow.  The following cut identities keep both
effects explicit.
Give a row vertex sign $\sigma_v=1$ and a column vertex sign
$\sigma_v=-1$.  For a small-cell incidence $(v,s)$, let
\[
a_{v,s}=\begin{cases}
x_s,&v\text{ is the row of }s,\\
b_s-y_s,&v\text{ is the column of }s,
\end{cases}
\qquad
\Delta a_{v,s}=a_{v,s}-X^*_s.
\]
For a large cell write $\Delta z_e=z_e-X^*_e$, including solved
tree cells in this notation.  Remove a tree edge $t$, and let $U_t$
be the vertex set on the side away from the root; let $v_t$ be the
endpoint of $t$ in $U_t$.  Sum the differences of the nonroot margin
equations over $U_t$, with signs $\sigma_v$.  Internal large edges
cancel.  If $v(e,U_t)$ is the unique endpoint in $U_t$ of a free edge
crossing the cut, the result is
\begin{equation}\label{eq:forest-cut}
\sigma_{v_t}\Delta z_t
=-\sum_{\substack{e\in\mathcal F\\e\text{ crosses }U_t}}
\sigma_{v(e,U_t)}\Delta z_e
-\sum_{\substack{s\in\mathcal S,\ v\in U_t\\v\text{ incident to }s}}
\sigma_v\Delta a_{v,s}.
\end{equation}
The last sum is over incidences, including both incidences when a
small edge has both endpoints in $U_t$; its two values can differ
for an unbalanced profile.

Each crossing free edge has $t$ on its tree path.  Hence
$|\Delta z_e|\le b_e+1\le2b_t$ even on the extended real free box.
Every small incidence satisfies
$|\Delta a_{v,s}|\le b_s<W\le b_t$.  There are at most $d$ free
edges and $2d$ small incidences.  Equation~\eqref{eq:forest-cut}
therefore gives $|\Delta z_t|\le4db_t$.

For a whole large-cell component $U$ with root $v_0$, the analogous
signed sum cancels every large-cell contribution and gives
\[
\sigma_{v_0}\delta_{v_0}
=\sum_{\substack{s\in\mathcal S,\ v\in U\\v\text{ incident to }s}}
\sigma_v\Delta a_{v,s}.
\]
This proves independence from all free values, also for a singleton
component.  Summing absolute values over components yields
$\sum_{v_0}|\delta_{v_0}|\le2\sum_{s\in\mathcal S}b_s\le2dW$.
Since $X^*_t\in[0,b_t]$, the displacement bound now implies
\[
0\le D\le A(2dW+4d^2)\le10Ad^2(W+1)=H.
\]
It also places $z_t$ in
$[-4db_t,(4d+1)b_t]\subseteq[-10db_t,10db_t]$.

Finally, subtract~\eqref{eq:forest-cut} for two free vectors with the
same small profile.  Each coefficient in normalized coordinates is
$\pm(b_e+1)$ for an edge crossing the cut, and its absolute value
divided by $b_t$ is at most two.  Thus
\[
\frac{|z_t(u)-z_t(u')|}{b_t}
\le2d\|u-u'\|_\infty.
\]
Distance from an interval is convex and is 1-Lipschitz.  Composing
with the affine tree values and summing at most $d$ tree terms proves
convexity and~\eqref{eq:forest-lipschitz}; the root terms are constant
in the free coordinates.
\end{proof}

Define the small-profile domain
\[
\mathcal P=\left\{(x,y)\in\prod_{s\in\mathcal S}
\{0,\ldots,b_s\}^2:
\sum_{s\in\mathcal S}(x_s+y_s)=p\right\}.
\]
For each integer $j\in\{0,\ldots,H\}$, set
\begin{equation}\label{eq:hj}
h_j(x,y)=
\sum_{z\in\prod_{e\in\mathcal F}\{0,\ldots,b_e\}}
2^{-jD(x,y,z)/H},\qquad (x,y)\in\mathcal P.
\end{equation}
When needed, extend $h_j$ by zero outside $\mathcal P$.  On
$\mathcal P$ it is strictly positive, including at unbalanced
profiles.  With the empty-product convention, put
$M_f=\prod_{e\in\mathcal F}(b_e+1)$.  Lemma~\ref{lem:forest} gives
\[
M_f2^{-j}\le h_j(x,y)\le M_f.
\]
An empty free block contributes one term, of value $2^{-jD(x,y)/H}$.
If there are no small cells, then $p=0$ and $\mathcal P$ contains
only the empty profile.

The partition functions of interest are
\begin{equation}\label{eq:Cj}
C(j)=\sum_{\substack{(x,y)\in\mathcal P\\
x_s+y_s=b_s\ (s\in\mathcal S)}}h_j(x,y).
\end{equation}
In particular,
$C(0)=M_f\prod_{s\in\mathcal S}(b_s+1)$ is explicitly computable.
Equations~\eqref{eq:hj} and~\eqref{eq:Cj} define finite mathematical
sums.  Section~\ref{sec:evaluation} and the subsequent annealing
algorithm estimate them using polynomially many operations on
binary representations, without enumerating the free ranges or the
small-profile domain. The outer counting algorithm in
Section~\ref{sec:algorithm} estimates the adjacent ratios in the identity
\[
 C(H)=C(0)\prod_{j=0}^{H-1}\frac{C(j+1)}{C(j)}.
\]

\begin{proposition}[Contamination bound]\label{prop:contamination}
For the feasible tightened instance,
\[
Z\le C(H)\le(1+\epsilon/32)Z.
\]
\end{proposition}

The lower bound counts each feasible table with weight one.  For the
upper bound, strong penalties immediately control root violations and
large tree violations; tree values just outside their intervals require
a marginal estimate.  We prove that estimate and complete the proof of
Proposition~\ref{prop:contamination} in
Section~\ref{sec:integer-signatures}.
